\usepackage{amssymb}
\usepackage{amsthm}
\usepackage{fancyhdr}
\usepackage{xcolor}
\usepackage{tikz}
\usetikzlibrary{graphs,graphdrawing,fit,backgrounds}
\usegdlibrary{force}

% A5 pages are too narrow for robust running section titles, especially
% with long mathematical headings. Use quiet page-number headers instead.
\pagestyle{fancy}
\fancyhf{}
\fancyhead[LE,RO]{\thepage}
\renewcommand{\headrulewidth}{0pt}
\fancypagestyle{plain}{%
  \fancyhf{}%
  \fancyfoot[C]{\thepage}%
  \renewcommand{\headrulewidth}{0pt}%
}
\emergencystretch=2em

% Emit an index hyperlink target only once, even if the same term or notation
% is defined in more than one indexed place.
\makeatletter
\newcommand{\indexhypertarget}[1]{%
  \@ifundefined{indexanchor@#1}{%
    \global\@namedef{indexanchor@#1}{}%
    \hypertarget{#1}{}%
  }{}%
}
\makeatother
\newcommand{\notationindexpage}[2]{\indexhypertarget{#1}\hyperpage{#2}}
\newcommand{\termindexpage}[2]{\indexhypertarget{#1}\hyperpage{#2}}
\makeatletter
\@ifundefined{ifbooklinkhighlight}{%
  \newif\ifbooklinkhighlight
  \booklinkhighlightfalse
}{}
% \booklinkanchorhere places a named destination at the *exact pen baseline*. A
% bare \hypertarget at a paragraph/list start runs in vertical mode, where the
% destination node lands in the inter-line space ~one line (about 16pt) above the
% first line's baseline. That offset exceeds the line spacing (~12pt), so a
% vertical-mode start destination sits above the *previous* line's baseline and the
% viewer cannot tell which line the span begins on. \leavevmode forces horizontal
% mode first (firing \everypar, so indentation is unaffected), so the destination
% attaches to the first line's baseline — the true span boundary. A bare
% \booklinkanchor is for ends, which the filter injects right after the last
% content character (already horizontal mode, already exact); forcing \leavevmode
% there could open a stray paragraph.
\newcommand{\booklinkanchor}[1]{\hypertarget{#1}{}}
\newcommand{\booklinkanchorhere}[1]{\leavevmode\hypertarget{#1}{}}
\makeatother
\ExplSyntaxOn
\int_new:N \l__booklink_entry_int
\keys_define:nn { booklink }
  {
    entry .int_set:N = \l__booklink_entry_int,
    entry .initial:n = 0,
    kind .code:n = {}
  }
% The booklink macros drop zero-width hypertargets (booklink-entry-N and
% booklink-entry-N-end) marking each span's start and end, so the formalization
% viewer can resolve the span's PDF extent and draw its overlay. They no longer
% color the text itself: the viewer's per-kind overlay is the single source of
% booklink coloring, so the rendered PDF stays visually clean. The anchors are
% un-raised (see \booklinkanchor above) so their coordinates are the exact span
% boundaries.
% At a list-item start the span begins with \item; \leavevmode *before* \item adds
% horizontal material ahead of it and breaks the list ("missing \item"). So when an
% \item is next, let it run first, then force horizontal mode (safe now) for the
% exact item-baseline anchor. Otherwise force horizontal mode directly so a
% paragraph/vertical-mode start lands on its first line's baseline.
\NewDocumentCommand{\BooklinkStart}{O{}m}
  {
    \keys_set:nn { booklink } { entry=0, #1 }
    \peek_meaning_remove:NTF \item
      { \item \exp_args:Nx \booklinkanchorhere { booklink-entry-\int_use:N \l__booklink_entry_int } }
      { \exp_args:Nx \booklinkanchorhere { booklink-entry-\int_use:N \l__booklink_entry_int } }
  }
\NewDocumentCommand{\BooklinkEnd}{O{}m}
  {
    \keys_set:nn { booklink } { entry=0, #1 }
    \exp_args:Nx \booklinkanchor { booklink-entry-\int_use:N \l__booklink_entry_int -end }
  }
\ExplSyntaxOff

% Formalization-skip anchors. The book filter wraps each `formalization: skip`
% region with \SkipStart{key}/\SkipEnd{key}; these drop zero-width hypertargets
% (skip-key-start / skip-key-end) so the formalization viewer can resolve the
% region's PDF coordinates and draw an overlay band. Gated on the same debug
% flag as the Booklink coloring, so the released book emits nothing and stays
% unchanged; only the booklink-highlight (debug) PDF carries the anchors.
\newcommand{\SkipStart}[1]{\ifbooklinkhighlight\hypertarget{skip-#1-start}{}\fi}
\newcommand{\SkipEnd}[1]{\ifbooklinkhighlight\hypertarget{skip-#1-end}{}\fi}

% Theorem/proof structure. Use these environments for genuine mathematical
% statements and proofs instead of raw Markdown bold labels such as
% **Theorem.**, **Lemma.**, **Corollary.**, **Claim.**, or **Proof.**.
% They are unnumbered for now so source cleanup does not change the book's
% visible numbering scheme.
\theoremstyle{plain}
\newtheorem*{theorem*}{Theorem}
\newtheorem*{lemma*}{Lemma}
\newtheorem*{proposition*}{Proposition}
\newtheorem*{corollary*}{Corollary}
\newtheorem*{claim*}{Claim}
\newtheorem*{fact*}{Fact}
\theoremstyle{definition}
\newtheorem*{recall*}{Recall}
\newtheorem*{definition*}{Definition}
\newtheorem*{example*}{Example}
\newtheorem*{construction*}{Construction}
\newtheorem*{remark*}{Remark}

% Generic amsthm environment for statement headings that do not fit one of the
% named theorem-like environments above.
\newtheorem*{statement*}{Statement}

% Boldface Borel/projective-hierarchy classes
\newcommand{\bSigma}{{\symbf{\Sigma}}}
\newcommand{\bPi}{{\symbf{\Pi}}}
\newcommand{\bDelta}{{\symbf{\Delta}}}
\newcommand{\DeltaZero}{\Delta_0}

% Classical low-level Borel class names.
% notation-index: $\Fsigma$
\newcommand{\Fsigma}{F_\sigma}
% notation-index: $\Gdelta$
\newcommand{\Gdelta}{G_\delta}
% notation-index: $\GdeltaSigma$
\newcommand{\GdeltaSigma}{G_{\delta\sigma}}
% notation-index: $\FsigmaDelta$
\newcommand{\FsigmaDelta}{F_{\sigma\delta}}

% Lightface Borel/projective-hierarchy classes
\newcommand{\lSigma}{\Sigma}
\newcommand{\lPi}{\Pi}
\newcommand{\lDelta}{\Delta}

% The set of natural numbers
% notation-index: $\NN$
% lint-raw-notation: enforce
\newcommand{\NN}{{\mathbb{N}}}
% The set of integers
% notation-index: $\ZZ$
% lint-raw-notation: enforce
\newcommand{\ZZ}{{\mathbb{Z}}}
% The set of rational numbers
% notation-index: $\QQ$
% lint-raw-notation: enforce
\newcommand{\QQ}{{\mathbb{Q}}}
% The set of real numbers
% notation-index: $\R$
% lint-raw-notation: enforce
\newcommand{\R}{\mathbb{R}}
% The set of complex numbers
% lint-raw-notation: enforce
\newcommand{\C}{\mathbb{C}}

% Function-space notation: \funcspace{A}{B} = functions from A to B,
% rendered as B^A (right superscript). Switch to {}^{A}\!B by editing
% this single definition.
% notation-index: $\funcspace{A}{B}$
\newcommand{\funcspace}[2]{#2^{#1}}

% Power set of A.
% notation-index: $\PowerSet{A}$
\newcommand{\PowerSet}[1]{\mathcal P(#1)}

% Inclusion relation allowing equality.
% notation-index: $A \IncludedIn B$
% lint-raw-notation: enforce
\newcommand{\IncludedIn}{\mathrel{\subseteq}}
% Strict inclusion relation.
% notation-index: $A \StrictlyIncludedIn B$
% lint-raw-notation: enforce
\newcommand{\StrictlyIncludedIn}{\mathrel{\subsetneq}}
% Set difference A without B.
% notation-index: $\SetDifference{A}{B}$
% lint-raw-notation: enforce
\newcommand{\SetDifference}[2]{#1 \setminus #2}

% Named partial single-valued map.
% notation-index: $\PartialFunctionMap{f}{X}{Y}$
\newcommand{\PartialFunctionMap}[3]{#1 : \subseteq #2 \to #3}

% Relation composition: first R, then S, rendered as S \circ R.
% notation-index: $\RelationComposition{S}{R}$
\newcommand{\RelationComposition}[2]{#1\circ #2}
% Injective map notation.
% notation-index: $\InjectionMap{f}{X}{Y}$
\newcommand{\InjectionMap}[3]{#1 : #2 \hookrightarrow #3}
% Surjective map notation.
% notation-index: $\SurjectionMap{f}{X}{Y}$
\newcommand{\SurjectionMap}[3]{#1 : #2 \twoheadrightarrow #3}

% Semantic aliases for the main spaces used throughout the text.
% These preserve the existing rendering while giving the source a more
% descriptive vocabulary.
\newcommand{\BaireSpace}{\funcspace{\omega}{\omega}}
\newcommand{\CantorSpace}{\funcspace{\omega}{2}}
% notation-index: $\HilbertCube$
\newcommand{\HilbertCube}{\seqPower{[0,1]}{\omega}}
% notation-index: $\RealSequenceSpace$
\newcommand{\RealSequenceSpace}{\seqPower{\R}{\omega}}
% Infinite symmetric group on the natural numbers.
% notation-index: $\InfiniteSymmetricGroup$
\newcommand{\InfiniteSymmetricGroup}{S_\infty}

% Sequence powers are function-space notation with ordinal domain n,
% read as sequences of length n.
% notation-index: $A^n$
\newcommand{\seqPower}[2]{#1^{#2}}

% Word/block repetition: w repeated n times as a concatenated word.
% notation-index: $w^n$
\newcommand{\wordRepeat}[2]{#1^{#2}}

% Finite sequences with entries in A: \finseq{A} = A^{<\omega}.
% notation-index: $A^{<\omega}$
\newcommand{\finseq}[1]{#1^{<\omega}}

% Infinite subsets of A: \infiniteSubsets{A} = [A]^\omega.
% notation-index: $\infiniteSubsets{\omega}$
\newcommand{\infiniteSubsets}[1]{[#1]^\omega}

% Disjoint union: A \disjointUnion B.
% notation-index: $A\disjointUnion B$
\newcommand{\disjointUnion}{\mathbin{\sqcup}}

% External indexed disjoint union/topological sum.
% notation-index: $\ExternalDisjointUnion_{i\in I} A_i$
\newcommand{\ExternalDisjointUnion}{\mathop{\bigsqcup}\displaylimits}

% Internal binary disjoint union: A \internalDisjointUnion B.
% notation-index: $A\internalDisjointUnion B$
\newcommand{\internalDisjointUnion}{\mathbin{\uplus}}

% Internal indexed/family disjoint union: \InternalDisjointUnion_{i\in I} A_i.
% notation-index: $\InternalDisjointUnion_{i\in I} A_i$
\newcommand{\InternalDisjointUnion}{\mathop{\uplus}\displaylimits}

% Directed supremum in domain theory.
% notation-index: $\DirectedSup S$
\newcommand{\DirectedSup}{\mathop{\bigsqcup}\displaylimits}

% Restriction of a function or sequence: \restrict{f}{A}.
% notation-index: $\restrict{f}{A}$
\newcommand{\restrict}[2]{#1\mathord{\upharpoonright}_{#2}}

% Sequence concatenation: s \concat t for the concatenation of finite
% (or finite-and-infinite) sequences. Renders as a superscript frown
% (the descriptive-set-theory standard, per Kechris / Moschovakis).
% notation-index: $s\concat t$
\newcommand{\concat}{\mathbin{{}^\frown}}

% Initial-segment / prefix relation for finite and infinite sequences.
% notation-index: $s\prefixOf y$
\newcommand{\prefixOf}{\mathrel{\preceq}}
% Negated initial-segment / prefix relation.
% notation-index: $s\notPrefixOf y$
\newcommand{\notPrefixOf}{\mathrel{\not\preceq}}
% Proper initial-segment / proper-prefix relation.
% notation-index: $s\properPrefixOf y$
\newcommand{\properPrefixOf}{\mathrel{\prec}}
% Reverse initial-segment relation: t extends s.
% notation-index: $t\extendsPrefix s$
\newcommand{\extendsPrefix}{\mathrel{\succeq}}
% Restriction of an ambient set by a subscript condition.
% notation-index: $P_{\le t}$
\newcommand{\RestrictedSet}[2]{#1_{#2}}

% Cylinder above a finite sequence: \cyl{s} = N_s = set of infinite
% extensions of s in the ambient sequence space. Rendered as N_s
% (Kechris / Moschovakis convention; "N" for neighborhood). Kept as a
% macro so the notation is grep-able and a future rename touches one
% line. We avoid [s] because [T] is already taken for the body of a
% tree.
% notation-index: $\cyl{s}$
\newcommand{\cyl}[1]{N_{#1}}

% The body [T] of a tree T.
% notation-index: $\TreeBody{T}$
\newcommand{\TreeBody}[1]{[#1]}

% Open subsets of a topological space X.
% notation-index: $\OpenSets{X}$
\newcommand{\OpenSets}[1]{\mathcal O(#1)}
% Closed subsets of a topological space X.
% notation-index: $\ClosedSets{X}$
\newcommand{\ClosedSets}[1]{\mathcal F(#1)}
% Prefix up-sets of finite sequences over an alphabet A.
% notation-index: $\PrefixUpSets{A}$
\newcommand{\PrefixUpSets}[1]{\operatorname{Up}_{\prefixOf}(\finseq{#1})}
% Prefix subtrees of finite sequences over an alphabet A.
% notation-index: $\PrefixSubtrees{A}$
\newcommand{\PrefixSubtrees}[1]{\operatorname{Subtr}_{\prefixOf}(\finseq{#1})}
% Cylinder-code map from prefix up-sets to open sets.
% notation-index: $\CylinderCode{A}$
\newcommand{\CylinderCode}[1]{\Phi(#1)}
% Open-cylinder-code map from open sets to prefix up-sets.
% notation-index: $\OpenCylinderCode{U}$
\newcommand{\OpenCylinderCode}[1]{\Psi(#1)}
% Closed-tree-code map from closed sets to prefix trees.
% notation-index: $\ClosedTreeCode{F}$
\newcommand{\ClosedTreeCode}[1]{T_{#1}}

% Diameter of a subset in the ambient metric.
% notation-index: $\Diameter{A}$
\newcommand{\Diameter}[1]{\mathrm{diam}(#1)}
% Diameter of a metric subset: \MetricDiameter{d}{A} = sup_{x,y \in A} d(x, y).
% The first argument names the metric.
% notation-index: $\operatorname{diam}_{d}(A)$
\newcommand{\MetricDiameter}[2]{\operatorname{diam}_{#1}(#2)}
% Open metric ball with center x and radius r.
% notation-index: $\Ball{x}{r}$
\newcommand{\Ball}[2]{B(#1,#2)}
% Open metric ball with an explicitly named metric.
% notation-index: $\MetricBall{d}{x}{r}$
\newcommand{\MetricBall}[3]{B_{#1}(#2,#3)}
% Open metric ball in an explicitly named ambient space.
% notation-index: $\AmbientBall{X}{x}{r}$
\newcommand{\AmbientBall}[3]{B_{#1}(#2,#3)}
% Closure in the ambient topological space.
% notation-index: $\operatorname{cl}(A)$
\newcommand{\Closure}[1]{\operatorname{cl}(#1)}
% Closure in an explicitly named ambient space.
% notation-index: $\operatorname{cl}_X(A)$
\newcommand{\AmbientClosure}[2]{\operatorname{cl}_{#1}(#2)}
% Standard metric on sequence spaces such as 2^omega and omega^omega.
% notation-index: $d_{\mathrm{seq}}$
\newcommand{\SeqMetricName}{d_{\mathrm{seq}}}
% notation-index: $\SeqMetric{x}{y}$
\newcommand{\SeqMetric}[2]{\SeqMetricName(#1,#2)}
% Cylinder radius for the standard sequence metric.
% notation-index: $\CylinderBallRadius{k}$
\newcommand{\CylinderBallRadius}[1]{\rho_{#1}}

% Domain of a function.
% notation-index: $\DomainOf{f}$
\newcommand{\DomainOf}[1]{\mathrm{dom}(#1)}

% Cantor–Bendixson derivative.
% \cbderivative{A}         produces A'
% \cbderivative[\alpha]{F} produces F^{(\alpha)}
% Use-before-definition is enforced at filter time (see tools/book-filter.hs);
% \define{cbderivative} in a mathmeta block marks the definition point.
% notation-index: $\cbderivative[\alpha]{F}$
\newcommand{\cbderivative}[2][]{\if\relax\detokenize{#1}\relax#2'\else#2^{(#1)}\fi}

% Cofinality of a limit ordinal
% notation-index: $\cf(\lambda)$
\newcommand{\cf}{\mathrm{cf}}

% Substitution: \subst{x}{t} renders as [x <- t], meaning "send x to t".
% Used in syntactic substitution into formulas: \varphi\subst{x}{t}.
\newcommand{\subst}[2]{[#1 \leftarrow #2]}
% notation-index: $t_0\subst{x}{t}$, first-order terms
\newcommand{\SubstTerm}[3]{#1\subst{#2}{#3}}
% notation-index: $\varphi\subst{x}{t}$, first-order formulas
\newcommand{\SubstFormula}[3]{#1\subst{#2}{#3}}
% notation-index: $M\subst{x}{N}$, lambda terms
\newcommand{\SubstLambda}[3]{#1\subst{#2}{#3}}
% notation-index: $t\subst{x_1,\ldots,x_n}{a_1,\ldots,a_n}$, PCA terms
\newcommand{\SubstPCA}[3]{#1\subst{#2}{#3}}
% Free variables of a syntactic expression.
\newcommand{\FreeVars}[1]{\mathrm{FV}(#1)}
% notation-index: $\FreeVars{t}$, first-order terms
\newcommand{\FreeVarsTerm}[1]{\FreeVars{#1}}
% notation-index: $\FreeVars{\varphi}$, first-order formulas
\newcommand{\FreeVarsFormula}[1]{\FreeVars{#1}}
% notation-index: $\FreeVars{M}$, lambda terms
\newcommand{\FreeVarsLambda}[1]{\FreeVars{#1}}
% First-order quantifier binders, kept separate from the raw symbols \forall
% and \exists so mathmeta inference only sees actual quantified formulas.
\newcommand{\forallbind}[2]{\forall #1\,#2}
\newcommand{\existsbind}[2]{\exists #1\,#2}
% Lambda abstraction, kept separate from the raw symbol \lambda so mathmeta
% inference only sees actual abstraction syntax.
\newcommand{\lambdaab}[2]{\lambda #1.#2}

% Forcing relation: p \forces \varphi.
% notation-index: $p\forces\varphi$
\newcommand{\forces}{\Vdash}

% Gandy--Harrington topology.
% notation-index: $\GHTopology$
\newcommand{\GHTopology}{\tau_{\rm GH}}

% Category-theoretic notation.
% Category names.
% notation-index: $\SetCat$
\newcommand{\SetCat}{\mathbf{Set}}
% notation-index: $\TopCat$
\newcommand{\TopCat}{\mathbf{Top}}
% notation-index: $\PolCat$
\newcommand{\PolCat}{\mathbf{Pol}}
% notation-index: $\MesCat$
\newcommand{\MesCat}{\mathbf{Mes}}
% notation-index: $\DomOmegaCat$
\newcommand{\DomOmegaCat}{\mathbf{Dom}_\omega}
% notation-index: $\qPolAdmCat$
\newcommand{\qPolAdmCat}{\mathbf{qPol}_{\rm adm}}
% notation-index: $\EffCat$
\newcommand{\EffCat}{\mathbf{Eff}}
% notation-index: $\CatCat$
\newcommand{\CatCat}{\mathbf{Cat}}
% notation-index: $\AsmCat{A}$
\newcommand{\AsmCat}[1]{\mathbf{Asm}(#1)}
% notation-index: $\ModCat{A}$
\newcommand{\ModCat}[1]{\mathbf{Mod}(#1)}
% notation-index: $\RTCat{A}$
\newcommand{\RTCat}[1]{\mathbf{RT}(#1)}
% notation-index: $\HeytAlgCat$
\newcommand{\HeytAlgCat}{\mathbf{HeytAlg}}
% notation-index: $\ShCat{X}$
\newcommand{\ShCat}[1]{\mathbf{Sh}(#1)}
% notation-index: $\SetCat_U$
\newcommand{\SetCatU}{\mathbf{Set}_U}
% notation-index: $\Ob{\mathcal C}$
\newcommand{\Ob}[1]{\mathrm{Ob}(#1)}
% Hom-set in a category: morphisms A -> B in C.
% notation-index: $\HomSet{\mathcal C}{A}{B}$
\newcommand{\HomSet}[3]{#1(#2,#3)}
% notation-index: $\opcat{\mathcal C}$
\newcommand{\opcat}[1]{#1^{\mathrm{op}}}
% notation-index: $\functorCat{\mathcal C}{\mathcal D}$
\newcommand{\functorCat}[2]{#2^{#1}}
% notation-index: $\presheafCat{\mathcal C}$
\newcommand{\presheafCat}[1]{\SetCat^{\opcat{#1}}}
% notation-index: $\catExp{A}{B}$
\newcommand{\catExp}[2]{#2^{#1}}
% notation-index: $\evalMap$
\newcommand{\evalMap}{\mathrm{ev}}
% notation-index: $\adjunction{F}{U}$
\newcommand{\adjunction}[2]{#1 \dashv #2}
% notation-index: $\Subobjects{A}$
\newcommand{\Subobjects}[1]{\mathrm{Sub}(#1)}
% notation-index: $\SubobjectClassifier$
\newcommand{\SubobjectClassifier}{\Omega}

% First uncountable ordinal
% notation-index: $\omegaOneOrd$
% lint-raw-notation: enforce
\newcommand{\omegaOneOrd}{\omega_1}
% Church--Kleene ordinal.
% notation-index: $\churchKleeneOrd$
% lint-raw-notation: enforce
\newcommand{\churchKleeneOrd}{\omega_1^{CK}}
% First infinite cardinal (= cardinality of \NN)
% notation-index: $\alephNull$
% lint-raw-notation: enforce
\newcommand{\alephNull}{\aleph_0}
% Aleph cardinal indexed by an ordinal.
% notation-index: $\alephCard{\alpha}$
% lint-raw-notation: enforce
\newcommand{\alephCard}[1]{\aleph_{#1}}
% Beth cardinals.
% notation-index: $\bethCard{\alpha}$
% lint-raw-notation: enforce
\newcommand{\bethCard}[1]{\beth_{#1}}
% Countable infinitary logic.
% notation-index: $\LCountableInfinitaryLogic$
\newcommand{\LCountableInfinitaryLogic}{L_{\omega_1\omega}}
% Consistency of a theory.
% notation-index: $\ConsistentTheory{T}$
\newcommand{\ConsistentTheory}[1]{\operatorname{Con}(#1)}
% The AD/Wadge determinacy ordinal: supremum of ordinals which are
% surjective images of the reals.
% notation-index: $\ADTheta$
\newcommand{\ADTheta}{\Theta}

% Effective descriptive set theory.
% notation-index: $\BorelCodes$
\newcommand{\BorelCodes}{\mathrm{BC}}
% notation-index: $\WellFoundedTrees$
\newcommand{\WellFoundedTrees}{\mathrm{WF}}
% notation-index: $\HyperarithmeticalSets$
\newcommand{\HyperarithmeticalSets}{\mathrm{HYP}}
% notation-index: $\KleeneO$
\newcommand{\KleeneO}{\mathcal{O}}
